\metatitle{Dot action and Hessenberg varieties}
\metadescription{A computational and representation-theoretic guide to
Tymoczko's dot action on regular semisimple Hessenberg varieties, its GKM
description, and its connection with chromatic quasisymmetric functions.}


\section[dotAction]{The dot action}

The \defin{dot action}\label{tymoczkoDotAction} is a representation of $\symS_n$ on the cohomology
of certain \hyperref[hessenbergVarieties]{Hessenberg varieties}.
It was introduced by \name{Julianna Tymoczko} in the setting of equivariant
cohomology \cite{Tymoczko2008,Tymoczko2008b}.
For regular semisimple Hessenberg varieties, this action gives the
representation-theoretic side of the Shareshian--Wachs conjecture.
The conjecture was proved by \name{Patrick Brosnan} and \name{Timothy Chow}
\cite{BrosnanChow2018}, and independently by \name{Mathieu Guay-Paquet}
\cite{GuayPaquet2016x}.

The broader chromatic symmetric function background is discussed on the
\hyperref[chromaticQuasisymmetricHessenberg]{chromatic quasisymmetric
function page}.  Here we isolate the representation-theoretic mechanism:
start with a Hessenberg sequence, build the GKM (\name{Mark Goresky},
\name{Robert Kottwitz}, and \name{Robert MacPherson}) model for equivariant
cohomology, descend the dot action to ordinary cohomology, and finally convert
the resulting graded characters into a
\hyperref[frobeniusCharacteristic]{Frobenius characteristic}.


\subsection[dotActionHessenbergSequence]{Hessenberg sequences and varieties}

\label{hessenbergVarieties}
A \defin{Hessenberg sequence}\label{hessenbergSequence} is a weakly increasing sequence
$m=(m_1,\dotsc,m_n)$ such that
\[
i\leq m_i\leq n
\qquad\text{for all } i.
\]
Equivalently, it is a \defin{Hessenberg function}\label{hessenbergFunction}
$m:[n]\to[n]$ with
$m(i)\geq i$.  This is the same data that appears in the area-sequence
description of \hyperref[chromaticQuasisymmetricUnitIntervalGraph]{unit interval graphs},
although the indexing convention depends on whether one starts from the
Hessenberg function or from the Dyck-diagram picture.  In the diagram convention
used on the chromatic quasisymmetric function page, the Hessenberg sequence
records the number of non-square boxes in each row of the extended
$n\times n$ diagram; for instance the area sequence
$(0,1,2,1,1)$ corresponds to $m=(2,3,5,5,5)$.  We write $\avec$ for the
corresponding area sequence and $\Gamma_\avec$ for the associated graph.

Let $F(n)$ be the set of all flags
\[
F_1\subset F_2\subset \dotsb \subset F_n=\setC^n,
\qquad
\dim F_i=i.
\]
Let $D$ be a diagonal $n\times n$ matrix with distinct diagonal entries.
The \defin{type $A$ regular semisimple Hessenberg variety}
\label{regularSemisimpleHessenbergVariety} associated with $m$ is
\[
H(m)
\coloneqq
\{F\in F(n):D F_i\subseteq F_{m_i}\text{ for all }i\}.
\]
For background on this construction, see the thesis of \name{Nicholas Teff}
\cite{Teff2013} and \name{Julianna Tymoczko}'s introduction to equivariant
cohomology and homology \cite{Tymoczko2005}.

The same sequence also determines a natural unit interval graph
$\Gamma_\avec$ on vertex set $[n]$.  In the standard Hessenberg-function
convention, we place an edge $\{i,j\}$, with $i \lt j$, whenever
\[
j\leq m_i.
\]
For this graph, the Shareshian--Wachs chromatic quasisymmetric function is
\[
\chrom_{\Gamma_\avec}(\xvec;q)
 =
\sum_{\substack{\kappa:[n]\to\setP\\ \kappa \text{ proper}}}
x_{\kappa(1)}\dotsm x_{\kappa(n)}q^{\asc(\kappa)},
\]
where the orientation is the increasing orientation $i\to j$ for $i \lt j$.

The original Shareshian--Wachs conjecture \cite[Conj.~5.3]{ShareshianWachs2012}
predicts that the cohomology of $H(m)$ carries a graded $\symS_n$-action whose
Frobenius characteristic is
$\omega\chrom_{\Gamma_\avec}(\xvec;q)$.  This was proved by
\name{Patrick Brosnan} and \name{Timothy Chow} \cite{BrosnanChow2018}, and
independently by \name{Mathieu Guay-Paquet} \cite{GuayPaquet2016x}.
There is also a closely related regular nilpotent story involving acyclic
orientations and Tymoczko cells; see \cite{NovelliThibon2024x}.


\subsection[dotActionGKMModel]{The GKM model}

Let $T=(\setC^\ast)^n$ be the diagonal torus.  We write
\[
H_T^\ast(\mathrm{pt})=\setC[t_1,\dotsc,t_n].
\]
The $T$-fixed points of $H(m)$ are naturally indexed by
permutations $w\in\symS_n$.
The moment graph records the zero- and one-dimensional $T$-orbits in $H(m)$.
The corresponding moment graph, denoted $MG(\avec)$ on
\hyperref[chromaticQuasisymmetricMomentGraph]{the chromatic quasisymmetric
function page}, has these permutations as vertices.
In the right-multiplication convention used below, the edges are
\[
w\longleftrightarrow w s_{ij}
\qquad
\text{for }\{i,j\}\in E(\Gamma_\avec).
\]

The \defin{GKM description}\label{gkmDescription} realizes the equivariant
cohomology as a ring of tuples
\[
p=(p_w)_{w\in\symS_n},
\qquad
p_w\in \setC[t_1,\dotsc,t_n],
\]
satisfying one divisibility condition for each edge of the moment graph.
For every edge $\{i,j\}$ of $\Gamma_\avec$ and every $w\in\symS_n$, we require
\[
p_w-p_{w s_{ij}}
\quad\text{is divisible by}\quad
t_{w(i)}-t_{w(j)}.
\]
Here $s_{ij}$ is the transposition swapping $i$ and $j$.
Thus
\[
H_T^\ast(H(m))
=
\left\{
(p_w)_{w\in\symS_n}
:
p_w-p_{w s_{ij}}\in \langle t_{w(i)}-t_{w(j)} \rangle
\text{ for all } w \text{ and } \{i,j\}\in E(\Gamma_\avec)
\right\}.
\]

The ordinary cohomology is obtained by setting the equivariant parameters to
zero:
\[
H^\ast(H(m))
\cong
H_T^\ast(H(m))
\otimes_{\setC[t_1,\dotsc,t_n]}\setC,
\]
where $t_1,\dotsc,t_n$ act as zero on $\setC$.
Equivalently, we quotient by the ideal generated by the positive-degree
polynomial variables.


\subsection[dotActionDefinition]{Definition of the dot action}

The group $\symS_n$ acts on $\setC[t_1,\dotsc,t_n]$ by permuting the indices:
\(\sigma(t_i)\coloneqq t_{\sigma(i)}\).
Tymoczko's \defin{dot action on a GKM tuple}
\label{dotActionOnGKMTuple} is
\[
(\sigma\cdot p)_w
\coloneqq
\sigma\left(p_{\sigma^{-1}w}\right).
\]
The subscript $\sigma^{-1}w$ moves the fixed-point label, and the outer
$\sigma$ permutes the polynomial variables.  The GKM divisibility relations are
preserved, so this gives an action on equivariant cohomology.  Since the ideal
$(t_1,\dotsc,t_n)$ is $\symS_n$-stable, the action descends to ordinary
cohomology.

The grading convention used in the chromatic formula is that polynomial degree
$d$ corresponds to cohomological degree $2d$.


\subsection[dotActionFrobeniusPipeline]{From the action to a symmetric function}

We now describe the steps needed to do computations with Tymoczko's dot action.
\begin{enumerate}
\item
Fix a Hessenberg sequence $m$, form the corresponding area sequence $\avec$,
and form the graph $\Gamma_\avec$.

\item
For each degree $d$, consider the vector space of all assignments
\[
\symS_n \to \setC[t_1,\dotsc,t_n]_d.
\]
This vector space has dimension
\[
n!\binom{n+d-1}{n-1},
\]
since each of the $n!$ tuple entries carries an independent copy of
\(\setC[t_1,\dotsc,t_n]_d\).

\item
Translate each GKM divisibility relation into linear equations in the
coefficients of the homogeneous polynomials;
the edges of $\Gamma_\avec$ determine which equations occur.
The solution space $M_d$ is the degree $d$ part of the equivariant GKM module.


\item
To pass to ordinary cohomology in degree $d$, quotient $M_d$ by
\[
R_d=\sum_{i=1}^n t_i M_{d-1}.
\]
Equivalently, $R_d$ is the degree $d$ part of
$(t_1,\dotsc,t_n)H_T^\ast(H(m))$.

\item
For each partition $\mu\vdash n$, choose a representative $\sigma_\mu$ of the
corresponding conjugacy class and compute the trace
\[
\chi_d(\mu)
=
\mathrm{tr}\left(
\sigma_\mu:
H^{2d}(H(m))\to H^{2d}(H(m))
\right).
\]
Since characters are class functions, the trace only needs to be computed on
one representative of each conjugacy class.
In practice, choose a basis for the quotient $M_d/R_d$, compute the matrix of
$\sigma_\mu$ on this quotient, and take its trace.  The trace is independent
of the chosen quotient basis.

\item
These trace values determine the
\hyperref[computingSnCharacters]{Frobenius characteristic}:
\[
\frobChar_d
=
\sum_{\mu\vdash n}\chi_d(\mu)\frac{\powerSum_\mu}{z_\mu}.
\]
The full graded Frobenius characteristic is
\[
\sum_d q^d \frobChar_d.
\]
\end{enumerate}


\subsection[dotActionShareshianWachsTheorem]{The Shareshian--Wachs theorem}

The main theorem connecting the above representation to symmetric functions is
the following.

\begin{theorem}[Brosnan--Chow, Guay-Paquet]
Let $m$ be a Hessenberg sequence, and let $\Gamma_\avec$ be the corresponding
natural unit interval graph.  Then
\[
\sum_{d\geq 0}
q^d
\frobChar\left(H^{2d}(H(m))\right)
=
\omega \chrom_{\Gamma_\avec}(\xvec;q).
\]
\end{theorem}

The right-hand side is the \hyperref[chromaticQuasisymmetric]{chromatic
quasisymmetric function} with the
\hyperref[standardInvolution]{$\omega$ involution} applied.  Consequently,
the theorem implies \hyperref[schurPositivity]{Schur positivity} of
$\omega\chrom_{\Gamma_\avec}(\xvec;q)$.

\name{Mathieu Guay-Paquet} gives a representation-level explanation of the
modular relation using divided difference operators on the GKM model
\cite{GuayPaquet2025x}.  In the stable and almost-stable cases these
operators decompose the equivariant cohomology module into dot-action
subrepresentations, categorifying the modular relation satisfied by
\hyperref[chromaticQuasisymmetric]{chromatic quasisymmetric functions}
\cite[Thm.~5.1 and
Thm.~5.3]{GuayPaquet2025x}.

\name{Tatsuya Horiguchi}, \name{Mikiya Masuda}, and \name{Takashi Sato}
give elementary GKM-theoretic proofs of the modular law for the geometric
objects behind the Shareshian--Wachs and unicellular LLT correspondences
\cite{HoriguchiMasudaSato2023x}.  This includes the regular semisimple
Hessenberg varieties and the twin varieties appearing in the LLT setting.

\name{Syu Kato} gives a different geometric realization of the chromatic
symmetric function of a unit interval graph using a smooth projective variety
$X_\Gamma$ and a map to the affine Grassmannian
\cite{Kato2024x}.  The decomposition theorem and geometric Satake identify
graded multiplicity spaces whose Schur-character sum is the chromatic
symmetric function \cite[Thm.~A]{Kato2024x}.  Kato also formulates a
geometric Stanley--Stembridge conjecture which would imply the
Shareshian--Wachs conjecture for these varieties.

\name{Kyle Salois} studies the regular semisimple Hessenberg varieties with
$h=(h(1),n,\dotsc,n)$ and the transpose family
$h'=((n-1)^{n-m},n^m)$ \cite{Salois2026}.  For these families he constructs
higher \hyperref[spechtModules]{Specht} bases and permutation bases for the
cohomology rings, together with bijections to $P_h$-tableaux.  The
permutation-basis result gives a
geometric explanation of the corresponding elementary-positivity statement.


\subsection[dotActionTypeBC]{Types B and C}

\name{Nathan R. T. Lesnevich} extends the same GKM/spline
viewpoint to types $B$ and $C$ \cite{Lesnevich2025x}.  The vertices are the
signed permutations $\symB_n$, and an order ideal $H$ in the positive roots
determines the signed transpositions and edge relations.  The resulting spline
module $\mathcal M_H$ carries a dot action and has two natural quotients: a
left quotient giving the ordinary-cohomology representation for regular
semisimple Hessenberg varieties in types $B$ and $C$, and a right quotient
related to the corresponding isospectral-matrix manifolds and unicellular LLT
polynomials.  Lesnevich computes the degree-one pieces of both dot-action
representations for all such $H$ directly from the combinatorics of these signed
transpositions.
